\documentclass[a4paper,final]{amsart}
\usepackage{mathtools}\mathtoolsset{showonlyrefs,showmanualtags,mathic=true}
   \usepackage[stixtwo]{fontsetup}
\usepackage[dvipsnames]{xcolor}
\usepackage{tikz-cd}
  \tikzcdset{arrow style=tikz, diagrams={>=stealth}}
  \usetikzlibrary {arrows.meta}
\usepackage[marginratio=1:1,scale=.72]{geometry}
\usepackage{caption}
  \captionsetup{width={0.75\textwidth}}
\usepackage{enumitem}
   \setlist{leftmargin=0pt,itemindent=2em,labelsep=1em}
\usepackage[breaklinks=true,colorlinks=true,allcolors=black]{hyperref}
\usepackage{showlabels}
   \renewcommand{\showlabelfont}{\tt\color{magenta}}
   \showlabels{bib}
\usepackage{amsrefs}
    \renewcommand{\eprint}[1]{\href{https://arxiv.org/abs/#1}{arXiv:#1}}
\usepackage{microtype}
\usepackage{tabularray}

\allowdisplaybreaks 

\usepackage{amsthm}
 \theoremstyle{definition}
   \newtheorem{definition}{Definition}
   \newtheorem{example}[definition]{Example}
 \theoremstyle{plain}
   \newtheorem{theorem}[definition]{Theorem}
   \newtheorem{proposition}[definition]{Proposition}
   \newtheorem{corollary}[definition]{Corollary}
 \theoremstyle{remark}
   \newtheorem{remark}[definition]{Remark}

\renewcommand{\qedsymbol}{\(\QED\)}

\DeclareMathOperator{\Hom}{Hom}
\DeclareMathOperator{\Ker}{Ker}
\DeclareMathOperator{\Cone}{Cone}
\DeclareMathOperator{\rk}{rk}
\DeclareMathOperator{\crk}{crk}

\title{
Remarks on Poincar\'e duality of matroids
}
\author{
Ryosuke Murooka \\
}
\address{Graduate School of Mathematics, Nagoya University, Nagoya, Japan}
\email{\href{mailto:ryosuke.murooka.c1@math.nagoya-u.ac.jp}{ryosuke.murooka.c1@math.nagoya-u.ac.jp}}
\date{}
%%%%%%%%%%%%%%%%%%%%%%%%%%%%%%%%%%%%%%%%%%%%%%%
\begin{document}
\begin{abstract}
Adiprasito, Huh, and Katz show Poincar\'e duality of matroids using matroidal flips, which are purely combinatorial notions.
Later, Hsiao, Karu, and Yang investigate the Hodge theory of matroids via star subdivisions, which correspond to blow-ups of toric varieties.
This is an expository paper reviewing a simple proof of Poincar\'e duality of matroids using their method of star subdivisions and Keel's lemma from algebraic geometry.
\end{abstract}
\maketitle
\section{Introduction}
Let \(n, r\) be non-negative integers.
In this paper, a \emph{matroid} \(M\) refers to a loopless matroid of rank \(r+1\) on an \((n+1)\)-set \(E\).
Write \(\symcal{L}(M)\) for the lattice of flats of \(M\).
Let \(\bigoplus_{i \in E}\symbf{Z}e_i\) be the free abelian group with \(E\) as its basis.
Put \(e_S \coloneq\sum_{i \in S}e_i\) for a subset \(S\subset E\), \(N_E\coloneq \bigl(\bigoplus_{i\in E}\symbf{Z}e_i\bigr)/\symbf{Z}e_E\), and \(N_{E,\symbf{R}}\coloneq N_E\otimes_{\symbf{Z}}\symbf{R}\).
The \emph{Bergman fan} \(\Sigma_M\) of a matroid \(M\) is the pure \(r\)-dimensional unimodular fan 
\begin{equation}
\big\{\,\Cone (e_F:F \in\symcal{F}) : \text{\(\symcal{F}\subset \symcal{P}(M)\) is a flag}\,\big\}
\end{equation}
in \(N_{E,\symbf{R}}\).
Here, \(\symcal{P}(M)\coloneq \symcal{L}(M)\setminus \{\emptyset,E\}\), and \(\symcal{F} \subset \symcal{P}(M)\) is called a \emph{flag} if it is linearly ordered with respect to the inclusion relation.
Associated to \(\Sigma_M\), we can consider the \(n\)-dimensional smooth quasi-projective complex toric variety \(X_{\Sigma_M}\), which is possibly non-compact.
Then, the \emph{Chow ring} \(A^{\bullet}(M)\) of \(M\) is defined to be the Chow ring \(A^{\bullet}\raisebox{-.25ex}{\big(}X_{\Sigma_M}\raisebox{-.25ex}{\big)}\) of \(X_{\Sigma_M}\).
Adiprasito, Huh, and Katz show Poincar\'e duality of \(M\) in the following sense.
\begin{theorem}[{\cite[Theorem 6.19]{AHK_18}}]\label{thm_2}
For a matroid \(M\) of rank \(r+1\), \(A^{\bullet}(M)\) satisfies Poincar\'e duality of dimension \(r\), i.e., it satisfies
\begin{enumerate}[label=$\arabic*)$]
\item there exists a group isomorphism \(\deg\colon A^r(M)\similarrightarrow \symbf{Z}\), called a \emph{degree map},
\item for every integer \(q > r\), we have \(A^q(M) = 0\), and 
\item for every integer \(q \leqq r\), the map
\begin{equation}
A^{r-q}(M) \to \Hom(A^q(M),\symbf{Z});\quad x \mapsto \big(y \mapsto\deg(xy)\big)
\end{equation}
is a group isomorphism, i.e., the multiplication \(A^q(M)\times A^{r-q}(M) \to A^r(M)\similarrightarrow \symbf{Z}\) is a perfect pairing.
\end{enumerate}
\end{theorem}
The key concepts of their proof are matroidal flips, which are purely combinatorial technical notions.
On the other hand, motivated by \cite{BHMPBW_22}, Hsiao, Karu, and Yang \cite{HKY_26} calculate mixed volumes of matroids using the deletion of matroids.
For their computation, they utilize the following important fact.
For a non-coloop \(a\in E\) of \(M\) and the deletion \(M\backslash a\), the projection \(\pi\colon N_E \to N_{E\setminus\{a\}}\) naturally induces a morphism of fans \(\pi\colon \Sigma_M \to \Sigma_{M\backslash a}\), and this morphism factors as
\begin{equation}\label{eq_1}\tag{$\clubsuit$}
\begin{tikzcd}[column sep =3em]
\Sigma_M = \Sigma_{m} \arrow[r,"\pi_{m}"] & \Sigma_{m-1} \arrow[r,"\pi_{m-1}"]& \cdots \arrow[r,"\pi_{2}"]& \Sigma_{1} \arrow[r,"\pi_{1}"]& \Sigma_0 \arrow[r,"\pi_0"]& \Sigma_{M\backslash a},
\end{tikzcd}
\end{equation}
where the morphism \(\pi_0\colon \Sigma_0\to \Sigma_{M\backslash a}\) is induced by the projection \(\pi\colon N_E \to N_{E\setminus\{a\}}\), \(\pi_i\colon \Sigma_{i} \to \Sigma_{i-1}\) is induced by the identity map, and \(\Sigma_i\) is a star subdivision of \(\Sigma_{i-1}\) at a 2-dimensional unimodular cone for each \(i=1,\dotsc,m\).
Using the factorization \eqref{eq_1}, Hsiao \cite{Hsiao_25} gives a simpler combinatorial proof that \(A^{\bullet}(M)\otimes_{\symbf{Z}}\symbf{R}\) satisfies the hard Lefschetz property, the Hodge-Riemann relations, and therefore Poincar\'e duality.
We note that Amini and Piquerez \cite[Section 1.9]{AP_23} also sketched a similar proof of these properties for matroids via star subdivisions, relying on the weak factorization theorem for tropical fans.
In \S \ref{sec_1}, we will give a proof of Theorem \ref{thm_2} in the language of algebraic geometry using \eqref{eq_1} and Keel's lemma.
%%%%%%%%%%%%%%%%%%%%%%%%%%
\section{Keel's lemma}\label{sec_2}
In this section, we recall an algebro-geometric lemma, which is known as Keel's lemma.
\begin{theorem}[Keel's lemma, cf.\ {\cite[Appendix, Theorem 1]{Keel_92}}]\label{thm_1}
Let \(X\) be a smooth variety, \(i\colon Z\hookrightarrow X\) a smooth closed subvariety of codimension \(z\), \(\pi\colon \widetilde{X}\to X\) the blow-up of \(X\) along \(Z\), and \(E\) the exceptional divisor, that is, we consider the pullback diagram
\begin{equation}
\begin{tikzcd}
E \arrow[r,hook]\arrow[d]\arrow[rd,phantom,"\pullback",very near start]& \widetilde{X}\arrow[d,"\pi"]\\
Z\arrow[r,hook,"i"'] & X.
\end{tikzcd}
\end{equation}
If the pullback map \(i^*\colon A^{\bullet}(X) \to A^{\bullet}(Z)\) is surjective, then the map \(\pi^*\colon A^{\bullet}(X)[T]\to A^{\bullet}(\widetilde{X});\ T\mapsto -[E]\) induces an isomorphism 
\begin{equation}
\begin{tikzcd}
\displaystyle \frac{A^{\bullet}(X)[T]}{(P(T), T\cdot \Ker i^*)} \arrow[r,"\text{\large$\sim$}" yshift=-.3ex]& A^{\bullet}(\widetilde{X})
\end{tikzcd}
\end{equation}
where \(P(T) \in A^{\bullet}(X)[T]\) is a monic polynomial of degree \(z\).
In particular, we have a decomposition
\begin{equation}\label{eq_2}\tag{$\diamondsuit$}
A^q(\widetilde{X}) = A^q(X)\oplus [E]\cdot A^{q-1}(Z) \oplus\dotsb\oplus [E]^{z-1}\cdot A^{q-z+1}(Z)
\end{equation}
for all \(q\), where we embed \(\pi^*\colon A^{\bullet}(X) \hookrightarrow A^{\bullet}(\widetilde{X})\), and for all positive integers \(p,q\), 
\begin{equation}
\begin{tikzcd}[column sep = 4em]
A^{q-p}(Z)\arrow[r,"(i^*)^{-1}"',"\text{\large$\sim$}" yshift=-.3ex]& A^{q-p}(X)/\Ker i^* \arrow[r,hook,"{[E]^p\cdot(-)}"']& A^{q}(\widetilde{X}).
\end{tikzcd}
\end{equation}
\end{theorem}
Now we go back to our toric case.
Let \(\Sigma\) be a unimodular fan in \(N_{E,\symbf{R}}\).
For a \(z\)-dimensional cone \(\sigma \in \Sigma\), let \(\Sigma_{\sigma}\) denote the star fan of \(\Sigma\) at \(\sigma\), and let \(\Sigma^*(\sigma)\) denote the star subdivision of \(\Sigma\) at \(\sigma\).
The orbit closure \(V(\sigma)\) associated to \(\sigma\) is the smooth toric subvariety \(X_{\Sigma_{\sigma}}\) of codimension \(z\), and \(\pi\colon X_{\Sigma^*(\sigma)}\to X_{\Sigma}\) is the blow-up of \(X_{\Sigma}\) along \(V(\sigma)\).
Since the ring \(A^{\bullet}(X_{\Sigma})\) is generated by the classes of torus-invariant divisors, the pullback homomorphism \(A^{\bullet}(X_{\Sigma}) \to A^{\bullet}\big(X_{\Sigma_{\sigma}}\big)\) is always surjective, so the assumption of Theorem \ref{thm_1} is satisfied if we substitute \(X=X_{\Sigma}\), \(Z=V(\sigma)=X_{\Sigma_{\sigma}}\), \(\widetilde{X}=X_{\Sigma^*(\sigma)}\).
We also note that the \(q\)-th degree part \(A^q(X_{\Sigma})\) is a finitely generated abelian group.
\begin{corollary}\label{cor_1}
We further assume that \(A^{\bullet}(X_{\Sigma})\) and \(A^{\bullet}\big(X_{\Sigma_{\sigma}}\big)\) satisfy Poincar\'e duality of dimension \(r\) and \(r-z\) respectively.
Then \(A^{\bullet}\big(X_{\Sigma^*(\sigma)}\big)\) satisfies Poincar\'e duality of dimension \(r\).
\end{corollary}
\begin{proof}
The decomposition \eqref{eq_2} tells us that \(A^r\big(X_{\Sigma^*(\sigma)}\big)\cong \symbf{Z}\) and \(A^q\big(X_{\Sigma^*(\sigma)}\big)=0\) for each \(q>r\).
By Poincar\'e duality of \(A^{\bullet}(X)\) and \(A^{\bullet}(Z)\), each summand of \eqref{eq_2} is a finitely generated free abelian group.
With an apropriately ordered basis chosen from each summand, the representation matrix of the bilinear pairing \(A^q\big(X_{\Sigma^*(\sigma)}\big)\times A^{r-q}\big(X_{\Sigma^*(\sigma)}\big)\to \symbf{Z}\) is a block lower triangular matrix 
\begin{equation}
\begin{bmatrix}
P_0 & O & \cdots &\cdots &O \\
O& P_1& \ddots& \ddots& \vdots \\
\vdots& \ast& \ddots & \ddots & \vdots\\
\vdots& \vdots& \ddots & \ddots & O\\
O&\ast&\ast&\ast&P_{z-1}
\end{bmatrix}.
\end{equation}
Poincar\'e duality of \(A^{\bullet}(X_{\Sigma})\) and \(A^{\bullet}(X_{\Sigma_{\sigma}})\) imply that \(\det P_q = \pm 1\) for all \(q=0,1,\dotsc,z-1\), and therefore the above representation matrix has determinant \(\pm 1\).
\end{proof}
%%%%%%%%%%%%%%%%%%%%%%%%%%%%%%%%%%%%%
\clearpage
\section{A proof of Poincar\'e duality by star subdivisions}\label{sec_1}
In this section, we prove Theorem \ref{thm_2} by induction on \(n\).  Before starting the proof, we recall the definition of a coloop.
\begin{definition}
An element \(a\in E\) is called a \emph{coloop} of \(M\) if the subset \(E\setminus\{a\}\) is a flat of \(M\), or equivalently, \(\dim \Sigma_{M\backslash a}< \dim \Sigma_M\).
\end{definition}
If all elements of \(E\) are coloops of \(M\), then \(M\) must be the Boolean matroid on \(E\).
The toric variety \(X_{\Sigma_M}\) is a smooth projective complex variety of dimension \(r=n\).
By the Jurkiewicz-Danilov theorem, the cycle map \(A^{\bullet}(X_{\Sigma_M})\to H^{2\bullet}(X_{\Sigma_M},\symbf{Z})\) is an isomorphism, and thus Poincar\'e duality of \(A^{\bullet}(M)\) follows in the usual sense.
In particular, this establishes the statement for \(n=0\).

We assume that \(n>0\) and there exists a non-coloop \(a\in E\) of \(M\).
If the singleton \(\{a\}\) is not a flat of \(M\), then the deletion map \(d\colon \symcal{L}(M) \to \symcal{L}(M\backslash a);\ F \mapsto F\setminus\{a\}\) is bijective, and \(\Sigma_M \subset (e_a-e_b)^{\perp}\) for some element \(b\) parallel to \(a\).
This implies that  \(X_{\Sigma_M}\cong X_{\Sigma_{M\backslash a}}\times \symbf{G}_{\symrm{m}}\), and hence \(A^{\bullet}(M)\cong A^{\bullet}(M\backslash a)\), where the right-hand side satisfies Poincar\'e duality of dimension \(r\) by the inductive hypothesis.
Hence we may assume that \(\{a\}\) is a flat of \(M\).
Now we construct the factorization \eqref{eq_1}.
Logically, the set \(\symcal{P}(M)\) can be partitioned into four disjoint subsets
\begin{align}
\symcal{P}(M) &=\big\{\,F:a \not\in F,\ F \cup\{a\}\in\symcal{P}(M)\,\big\}\sqcup\big\{\,F: a\not\in F,\ F \cup \{a\} \not\in \symcal{P}(M)\,\big\}\\
&\hspace{11em}\sqcup\big\{\,F: a\in F,\ F \setminus\{a\} \in \symcal{P}(M)\,\big\}\sqcup\big\{\,F: a\in F,\ F \setminus \{a\} \not\in \symcal{P}(M)\,\big\}\\
&\eqcolon \symcal{P}_1\sqcup\symcal{P}_2\sqcup\symcal{P}_3\sqcup\symcal{P}_4.
\end{align}
Under the deletion map \(d\colon \symcal{P}(M)\to \symcal{P}(M\backslash a)\), we have
\[
\symcal{P}(M\backslash a)=\symcal{P}_1\sqcup \symcal{P}_2 \sqcup d(\symcal{P}_4),\qquad d(\symcal{P}_3)\subset \symcal{P}_1,\qquad \symcal{P}(M/a)=d(\symcal{P}_3)\sqcup d(\symcal{P}_4).
\]
Consider the inverse map \(d^{-1}\colon \symcal{P}(M\backslash a)\to \symcal{P}(M)\setminus \symcal{P}_3\).
We construct a fan \(\Sigma_0\) in \(N_{E,\symbf{R}}\) by
\begin{align}
&\Sigma_0\coloneq \big\{\,\Cone(e_{d^{-1}(F)}:F \in \symcal{F}): \symcal{F}\underset{\text{flag}}{\subset}\symcal{P}(M\backslash a)\,\big\}\\
&\hspace{12em}\cup \big\{\,\Cone(e_{d^{-1}(F)}:F \in \symcal{F})+\Cone (e_a):\symcal{F}\underset{\text{flag}}{\subset} \symcal{P}(M/a)\,\big\}
\end{align}
(the latter cones are intended to recover the rays associated to \(\symcal{P}_3\)).
Observing that \(\Sigma_0\) has exactly one more ray than \(\Sigma_{M\backslash a}\) \big(the ray is \(\Cone(e_a)\)\big), one can show that the projection \(\pi_0\colon\Sigma_0 \to \Sigma_{M\backslash a}\) induces an isomorphism \(\pi_0^*\colon A^{\bullet}(M\backslash a) \similarrightarrow A^{\bullet}(X_{\Sigma_0})\).
We order the flats \(d(\symcal{P}_3)=\{F_1, F_2,\dotsc, F_m\}\) so that \(F_i\subset F_j\) implies that \(i \leqq j\).
Recall that, since \(d(\symcal{P}_3)\subset \symcal{P}_1\), we have \(d^{-1}(F_i)=F_i\) for all \(i\), and hence \(\Cone(e_a,e_{F_i})\in \Sigma_0\) for all \(i\).
Let \(\Sigma_1\) be the star subdivision of \(\Sigma_0\) at \(\Cone(e_{F_m},e_a)\).
\(\Sigma_1\) has an extra ray \(\Cone(e_{F_{m}\cup\{a\}})\), and all new cones arising from this subdivision are of the form
\begin{align}
\Cone (e_{d^{-1}(F)}: F \in \symcal{F})+\Cone (e_{F_m\cup\{a\}}) &\quad \text{where \(\symcal{F}\) is a flag of \(\symcal{P}(M/a)\), or}\\
\Cone (e_{d^{-1}(F)}: F \in \symcal{F})+\Cone(e_a)+\Cone (e_{F_m\cup\{a\}}) &\quad \text{where \(\symcal{F}\) is a flag of \(\symcal{P}(M/a)\setminus\{F_m\}\).}
\end{align}
The cones of the former form are already cones of \(\Sigma_M\).
By recursively constructing the star subdivision \(\Sigma_{i+1}\) of \(\Sigma_i\) at \(\Cone(e_{F_{m-i}},e_a)\) for each \(i=1,\dotsc,m-1\), we obtain \(\Sigma_{m}=\Sigma_M\) and the desired factorization 
\begin{equation}\tag{$\clubsuit$}
\begin{tikzcd}[column sep =3em]
\Sigma_M = \Sigma_{m} \arrow[r,"\pi_{m}"] & \Sigma_{m-1} \arrow[r,"\pi_{m-1}"]& \cdots \arrow[r,"\pi_{2}"]& \Sigma_{1} \arrow[r,"\pi_{1}"]& \Sigma_0 \arrow[r,"\pi_0"]& \Sigma_{M\backslash a}.
\end{tikzcd}
\end{equation}

Finally, the star fan of \(\Sigma_i\) at \(\Cone(e_{F_{m-i}},e_a)\) is isomorphic to the product \(\Sigma_{[\emptyset,F_{m-i}]}\times \Sigma_{[F_{m-i}\cup\{a\},E]}\), where \([\emptyset,F_{m-i}]\) denotes the restriction of \(M\) to \(F_{m-i}\) and \([F_{m-i}\cup\{a\},E]\) denotes the contraction of \(M\) by \(F_{m-i}\cup\{a\}\).
This implies that its Chow ring is isomorphic to \(A^{\bullet}([\emptyset,F_{m-i}])\otimes A^{\bullet}([F_{m-i}\cup\{a\},E])\), which satisfies Poincar\'e duality of dimension \((\rk_M(F_{m-i})-1)+(\crk_M(F_{m-i}\cup\{a\})-1)=r-2\) by the inductive hypothesis.
Applying Corollary \ref{cor_1} with \(z=2\), we conclude that \(A^{\bullet}(M)\) satisfies Poincar\'e duality of dimension \(r\).
\qed
%%%%%%%%%%%%%%%%%%
\clearpage

We give some remarks of the proof.
\begin{remark}
If we subdivide \(\Sigma_0\) at \(\Cone(e_{F_1},e_a)\), the resultant fan differs from \(\Sigma_M\).
For instance, when \(F_1\subset F_2\), \(\Cone(e_{F_1},e_{F_2},e_a)\) is divided into \(\Cone(e_{F_1\cup\{a\}},e_{F_2},e_a)\) and \(\Cone(e_{F_1\cup\{a\}},e_{F_1},e_{F_2})\).
The former has a face \(\Cone(e_{F_1\cup\{a\}},e_{F_2})\), which remains in the final fan and is not a cone of Bergman fan.
\end{remark}
\begin{remark}
In the language of tropical geometry, the fan \(\Sigma_0\) is known as the tropical modification of the tropical fan \(\Sigma_{M\backslash a}\) along the tropical divisor \(\Sigma_{M/a}\) with respect to a suitable piecewise linear function on \(\Sigma_{M\backslash a}\).
Then the isomorphism \(\pi^*\colon A^{\bullet}(M\backslash a)\to A^{\bullet}(X_{\Sigma_0})\) follows from \cite[Theorem 6.4]{AP_23}.
\end{remark}
\begin{remark}
Amini and Piquerez \cite[Theorem 7.9]{AP_23} show that \(X_{\Sigma_M}\) and \(X_{\Sigma_0}\) can be connected by a finite sequence of toric blowups and blowdowns (weak factorization for tropical fans).
Therefore, the sequence \(\Sigma_m \to \Sigma_{m-1} \to \dotsb \to \Sigma_0\) in \eqref{eq_1} provides a concrete example of this theorem for matroids.
Furthermore, it shows that \(X_{\Sigma_M}\) can be obtained from \(X_{\Sigma_0}\) by a finite sequence of blowups.
\end{remark}
%%%%%%%%%%%%%%%%%%%%%%%%%%%%%
\section*{Acknowledgements}
The author is grateful to his advisor, Yusuke Nakamura, for steady guidance and numerous insightful comments.
The author would like to express his sincere gratitude to Masahiko Yoshinaga for giving him helpful comments to improve the proof of Poincar\'e duality of matroids.
The author would also like to thank Kosuke Mizuno and Kentaro Tanaka for informative comments.
%%%%%%%%%%%%%%%%%%%%%%%%%%%%%
\begin{bibdiv}
\begin{biblist}
\bib{AHK_18}{article}{
   title={Hodge theory for combinatorial geometries},
   author={Adiprasito, Karim},
   author={Huh, June},
   author={Katz, Eric},
   journal={Ann.\ Math.\ (2)},
   date={2018},
   volume={188},
   number={2},
   pages={381--452},
   note={\href{https://doi.org/10.4007/annals.2018.188.2.1}{DOI 10.4007/annals.2018.188.2.1}},
}
\bib{AP_23}{article}{
  author={Amini, Omid},
  author={Piquerez, Matthieu},
  title={Hodge theory for tropical fans},
  date={2023},
  eprint={2310.15367v1},
  status={to appear in Bull. Soc. Math. France},
}
\bib{BHMPBW_22}{article}{
   title={A semi-small decomposition of the Chow ring of a matroid},
   author={Braden, Tom},
   author={Huh, June},
   author={Matherne, Jacob P.},
   author={Proudfoot, Nicholas},
   author={Wang, Botong},
   journal={Adv.\ Math.},
   date={2022},
   volume={409},
   pages={Article 108646},
   note={\href{https://doi.org/10.1016/j.aim.2022.108646}{DOI 10.1016/j.aim.2022.108646}},
}
\bib{Hsiao_25}{article}{
   title={Hodge theory of matroids by star subdivisions},
   author={Hsiao, Andy},
   date={2025},
   eprint={2509.10909v1},
   status={preprint},
}
\bib{HKY_26}{article}{
   title={Mixed Volumes of Matroids},
   author={Hsiao, Andy},
   author={Karu, Kalle},
   author={Yang, Jonathan},
   journal={Discrete Comput.\ Geom.},
   date={2026},
   note={\href{https://doi.org/10.1007/s00454-026-00832-y}{DOI 10.1007/s00454-026-00832-y}},
}
\bib{Keel_92}{article}{
   title={Intersection theory of moduli space of stable \(n\)-pointed curves of genus zero},
   author={Keel, Sean},
   date={1992},
   journal={Trans.\ Amer.\ Math.\ Soc.},
   volume={330},
   pages={545--574},
   note={\href{https://doi.org/10.1090/S0002-9947-1992-1034665-0}{DOI 10.1090/S0002-9947-1992-1034665-0}},
}
\end{biblist}
\end{bibdiv}
\end{document}
